\scanpage{48}
% archive sheet number: 47

\mspage{11}
mais sa quatrième fibre est

\[
\bar X_{(1,1)}=X_{/(1,1)}\simeq X,
\]

donc le diagramme \inserted{dans Cat} du type $\underline Q$ donnant naissance à $Y_{\underline Q}(X)$ n'est autre que

\diagram{\begin{tikzcd}[ampersand replacement=\&]
\& U\arrow[dr,no head] \&\\
U\cap U'\arrow[ur]\arrow[dr] \& \& X\\
\& U'\arrow[ur,no head] \&
\end{tikzcd}}

Mais ce diagramme \inserted{comm. dans Cat} n'est bien sûr pas quelc, sa propriété essentielle, pour les propriétés que nous avons en vue, me semble la suivante :

\diagram{\[
\left(\underbrace{\int
\begin{tikzcd}[ampersand replacement=\&,column sep=small,row sep=small]
\& U\\
\deleted{\unsure{U\cap U'|V}}\arrow[ur]\arrow[dr,no head] \&\\
\& U'
\end{tikzcd}}_{\substack{\unsure{\parallel|\Downarrow}\,\mathrm{déf}\\Y_{\underline\Psi}(X)\ \text{en tant}\\\text{qu'objet de Cat sans plus}}}\right)\longrightarrow X
\]}

\margin{(flèche définie pour tt diagramme comm.}

\diagram{\begin{tikzcd}[ampersand replacement=\&,column sep=small,row sep=small]
\& U\arrow[dr] \&\\
V\arrow[ur]\arrow[dr,no head] \& \& X\\
\& U'\arrow[ur] \&
\end{tikzcd}}

\margin{dans Cat}

est dans $W$. Et le fait général pour y satisfaire est le suivant

``\uline{Principe}'' $(?)$ Soit

\diagram{\begin{tikzcd}[ampersand replacement=\&,row sep=small]
\& X_a\\
X_0\arrow[ur]\arrow[dr] \&\\
\& X_b
\end{tikzcd}}

un diagramme dans Cat du type $\underline\Psi$, soit $\mathcal X$ la cat. cofibrée sur $\underline\Psi$ qu'il définit. Pour tout $\xi$
